% =============================================================================
\section{Extension formelle : fermeture SCHA en phase ferroélectrique}
\label{sec:core_ferroelectric_extension}
% =============================================================================

Cette annexe est une extension complémentaire. Elle n'est utilisée ni dans le
solveur numérique paraélectrique, ni dans les courbes, ni dans les prédictions
présentées dans ce rapport. Le calcul effectivement exécuté reste celui de la
branche centrée $\overline{\mathbf P}=0$ décrit dans
Sec.~\ref{ann:scha_reponse}.

Dans une éventuelle généralisation à la phase de symétrie brisée, le fond
statique homogène deviendrait une variable variationnelle :
\begin{equation}
\overline p_{\min}^{\alpha}(k,\omega_n)
=\sqrt{\beta\hbar}\,p_{\min}^{\alpha}(T)\frac{(2\pi)^3}{\Omega}
\delta^{(3)}(k)\delta_{n,0},
\qquad
p'^\alpha=p^\alpha-\overline p_{\min}^{\alpha}.
\label{eq:core_ferro_shift}
\end{equation}
La mesure gaussienne est centrée sur $p'=0$, tandis que le champ physique vaut
$\widehat p=p'+\overline p_{\min}$. L'énergie libre variationnelle dépendrait
alors simultanément du noyau anisotrope $A$ et des trois composantes du
centroïde :
\begin{equation}
F_{\mathrm{var}}(T;A,\mathbf p_{\min})
=-k_BT\ln Z_0[A]
+\left\langle V(\widehat{\mathbf p})-V_0(\mathbf p')\right\rangle_A.
\label{eq:core_ferro_variational_free_energy}
\end{equation}
L'évaluation de cette fonctionnelle ne requiert pas une nouvelle règle de
contraction. En posant $\widehat P_a=\overline P_a+\xi_a$, avec
$\langle\xi_a\rangle_A=0$ et $\langle\xi_a\xi_b\rangle_A=G_{ab}$, tout
moment est obtenu par l'identité compacte
\begin{equation}
\left\langle\prod_{j=1}^{r}\widehat P_{a_j}\right\rangle_A
=\sum_{\substack{S\subset\{1,\ldots,r\}\\r-|S|\;\mathrm{pair}}}
\left(\prod_{i\in S}\overline P_{a_i}\right)
\sum_{\text{appariements de }\bar S}
\prod_{(i,j)}G_{a_i a_j}.
\label{eq:core_ferro_general_wick_rule}
\end{equation}
Cette formule engendre mécaniquement les termes à deux, quatre, six et huit
jambes : tout terme comportant un nombre impair de fluctuations $\xi$ est nul.
Elle fournit donc une prescription calculable pour développer les vertices
locaux sans avoir à réécrire séparément leurs dizaines de contractions.

Pour rendre les équations à résoudre explicites, notons
$V_{\rm eff}[P]$ le potentiel complet après élimination des deux secteurs de
déformation. La stationnarité de la fonctionnelle précédente prend la forme
générale
\begin{equation}
\boxed{\quad A_{ab}(K)=A_{0,ab}(K)
+\left\langle
\frac{\delta^2V_{\rm eff}^{\rm anh}}
{\delta P_a(-K)\,\delta P_b(K)}
\right\rangle_{\overline P,G},
\qquad
0=\left\langle
\frac{\delta V_{\rm eff}}{\delta P_a}
\right\rangle_{\overline P,G}.\quad}
\label{eq:core_ferro_general_scha_equations}
\end{equation}
Dans les conventions intensives du présent rapport, la première de ces deux
conditions donne explicitement la fermeture ferroélectrique finale
\begin{mdframed}[linecolor=red!60!black,nobreak=true]
\small
\begin{equation}
\begin{aligned}
A_{ab}(K)={}&A_{0,ab}(K)
+3\beta_{ab\gamma\delta}
\big(G^{\gamma\delta}+\overline P^\gamma\overline P^\delta\big)
-\frac12\Lambda^{\rm hom}_{ab\gamma\delta}
\big(G^{\gamma\delta}+\overline P^\gamma\overline P^\delta\big)\\
&-\mathcal I^{\rm inh}_{ab}(K;A)
-\mathcal M^{\rm inh}_{a\gamma,b\delta}(K)
\overline P^\gamma\overline P^\delta\\
&+15\gamma_{ab\gamma\delta\epsilon\zeta}
\left[G^{\gamma\delta}G^{\epsilon\zeta}
+2\overline P^\gamma\overline P^\delta G^{\epsilon\zeta}
+\frac13\overline P^\gamma\overline P^\delta
\overline P^\epsilon\overline P^\zeta\right]\\
&+105\rho_{ab\gamma\delta\epsilon\zeta\eta\theta}
\Big[G^{\gamma\delta}G^{\epsilon\zeta}G^{\eta\theta}
+3\overline P^\gamma\overline P^\delta
G^{\epsilon\zeta}G^{\eta\theta}\\[-2pt]
&\hspace{33mm}
+\overline P^\gamma\overline P^\delta
\overline P^\epsilon\overline P^\zeta G^{\eta\theta}
+\frac1{15}\overline P^\gamma\overline P^\delta
\overline P^\epsilon\overline P^\zeta
\overline P^\eta\overline P^\theta\Big].
\end{aligned}
\label{eq:core_ferro_full_closure}
\end{equation}
\end{mdframed}
Cette écriture vaut pour les modes non uniformes $K\neq0$. Le terme éliminé
de déformation homogène,
$-\Lambda^{\rm hom}_{\alpha\beta\gamma\delta}
(\sum_iP_i^\alpha P_i^\beta)(\sum_jP_j^\gamma P_j^\delta)/(8N)$, induit en
outre une correction réservée au seul mode statique uniforme. Plutôt que de
l'insérer par des deltas dans le noyau général, on l'ajoute séparément : le
membre de droite de l'Eq.~\eqref{eq:core_ferro_full_closure}, évalué à $K=0$,
reçoit le terme supplémentaire
\begin{equation}
 -\Lambda^{\rm hom}_{a\gamma b\delta}
 \overline P^\gamma\overline P^\delta.
\label{eq:core_ferro_homogeneous_uniform_channel}
\end{equation}
Le terme
$\mathcal M^{\rm inh}_{a\gamma,b\delta}(K)
\overline P^\gamma\overline P^\delta$ est l'analogue provenant de la
déformation inhomogène, mais il conserve une dépendance générale en $K$.
Les deux doivent être conservés si la fermeture ferroélectrique est
effectivement implémentée. À $\overline P=0$, ils disparaissent et
l'Eq.~\eqref{eq:core_ferro_full_closure} redonne
l'Eq.~\eqref{eq:core_full_scha_tensor}.
\needspace{15\baselineskip}
La première équation détermine le noyau et la seconde le centroïde. Par
exemple, si l'on isole les vertices locaux du potentiel, elles contiennent
explicitement
\begin{equation}
\begin{aligned}
\Pi_{ab}^{\rm loc}={}&3\beta_{ab\gamma\delta}
\langle\widehat P_\gamma\widehat P_\delta\rangle
+5\gamma_{ab\gamma\delta\epsilon\zeta}
\langle\widehat P_\gamma\widehat P_\delta
\widehat P_\epsilon\widehat P_\zeta\rangle\\
&+7\rho_{ab\gamma\delta\epsilon\zeta\eta\theta}
\langle\widehat P_\gamma\widehat P_\delta
\widehat P_\epsilon\widehat P_\zeta
\widehat P_\eta\widehat P_\theta\rangle .
\end{aligned}
\label{eq:core_ferro_local_self_energy}
\end{equation}
\begin{equation}
\begin{aligned}
0={}&A_{0,ab}(0)\overline P_b
+\beta_{abcd}\langle\widehat P_b\widehat P_c\widehat P_d\rangle\\
&+\gamma_{abcdef}\langle\widehat P_b\widehat P_c\widehat P_d
\widehat P_e\widehat P_f\rangle
+\rho_{abcdefgh}\langle\widehat P_b\widehat P_c\widehat P_d
\widehat P_e\widehat P_f\widehat P_g\widehat P_h\rangle\\
&+\left\langle\frac{\delta V_{\rm strain}}{\delta P_a}\right\rangle .
\end{aligned}
\label{eq:core_ferro_centroid_equation}
\end{equation}
Les indices sont contractés selon la convention d'Einstein, et
$V_{\rm strain}$ désigne les termes homogène et inhomogène non locaux. Les
moments de ces deux équations se déduisent directement de
l'Eq.~\eqref{eq:core_ferro_general_wick_rule}; le système est donc défini sans
ambiguïté, même si sa résolution dépasse le périmètre numérique du présent
rapport.
Les deux conditions de stationnarité seraient
\begin{equation}
A_{ab}(K)=A_{0,ab}(K)
+\Pi_{ab}(K;A,\mathbf p_{\min}),
\qquad
\frac{\partial F_{\mathrm{var}}}
{\partial p_{\min}^{a}}=0.
\label{eq:core_ferro_stationarity}
\end{equation}
La self-énergie $\Pi$ contiendrait les contractions locales quartiques,
sextiques et octiques, ainsi que les secteurs $\Lambda^{\mathrm{hom}}$ et
$\mathcal M^{\mathrm{inh}}$, avec leurs termes croisés en
$p_{\min}^{a}p_{\min}^{b}$.

Les moments déplacés nécessaires à cette extension seraient générés par
\begin{equation}
\langle\widehat P_a\widehat P_b\rangle
=\overline P_a\overline P_b+G_{ab},
\label{eq:core_shifted_second_moment}
\end{equation}
et, pour le rang quatre,
\begin{equation}
\begin{aligned}
\langle\widehat P_a\widehat P_b\widehat P_c\widehat P_d\rangle
={}&\overline P_a\overline P_b\overline P_c\overline P_d\\
&+\overline P_a\overline P_bG_{cd}
+\overline P_a\overline P_cG_{bd}
+\overline P_a\overline P_dG_{bc}\\
&+\overline P_b\overline P_cG_{ad}
+\overline P_b\overline P_dG_{ac}
+\overline P_c\overline P_dG_{ab}\\
&+G_{ab}G_{cd}+G_{ac}G_{bd}+G_{ad}G_{bc}.
\end{aligned}
\label{eq:core_shifted_fourth_moment}
\end{equation}
Les rangs six et huit suivraient la même règle : choisir les facteurs portés
par le centroïde, puis apparier les fluctuations restantes.

Un solveur dédié à cette extension devrait traiter les neuf inconnues
\begin{equation}
\mathbf x=
\big(A_{xx},A_{yy},A_{zz},A_{xy},A_{xz},A_{yz},
\overline P_x,\overline P_y,\overline P_z\big),
\label{eq:core_ferro_unknown_vector}
\end{equation}
avec six résidus issus de $A-A_0-\Pi[A,\overline{\mathbf P}]$ et trois
conditions $\partial F_{\rm var}/\partial\overline P_a=0$. Ce solveur n'est pas
appelé dans la chaîne numérique auditée et aucune direction ferroélectrique
n'est utilisée pour produire les résultats du présent travail.
